\scanpage{4}
\mspage{- III-18 -}
$C_y$ ; il est immédiat que tout élément de cette algèbre peut s'écrir\unsure{e} \margin{\handwritten{$\ast$}}$d\handwritten{\otimes}(1/t)$, où $t\handwritten{\in}C$ \deleted{\illegible{1}} est tel que $t\handwritten{\notin}j_y\inserted{\handwritten{=0}}$, et $d\handwritten{\in}D$ ; elle, a par suite une base de la forme $(d_i\handwritten{\otimes}1)_{1\handwritten{\leq}i\handwritten{\leq}n}$, où $d_i\handwritten{\in}D$. Nous allons voir que l'on peut supposer\deleted{\illegible{1}} que les $d_i$, qui sont des sections de $\underline O_X$ au-dessus de $f^{-1}(U)$, sont des restrictions de sections de $\underline O_X$ \uline{au-dessus de} $X$. Remarquons pour cela que le préschéma $f^{-1}(y)$ a un espace de base formé d'un nombre fini de points $x_j$ ($1\handwritten{\leq}j\handwritten{\leq}m$ ; II,2,7,pro).17) formant un sous-espace \uline{discret} de $f^{-1}(U)$ ; l'algèbre \deleted{\illegible{1}} $D\handwritten{\otimes}_C C_y$, qui est l'anneau de ce schéma affine, est donc composée directe des sous-algèbres $\underline O_{x_j}$ (I,\uline{2,7,prop.19 et I,4,1,prop.4}), \margin{\handwritten{!}}et on peut \uline{évidemment} supposer que chacun des éléments $d_i\handwritten{\otimes}1$ est un élément de base
% slanted
\margin{\handwritten{? \unsure{contenu} dans un \unsure{des}}}de l'un des $\underline O_{x_j}$. Comme $X$ est affine, il y a pour chaque $j$ un élément $f_j\handwritten{\in\Gamma}(X,\underline O_X)=B$ tel que $x_j\handwritten{\in}X_{f_j}\handwritten{\subset}f^{-1}(U)$ \deleted{\unsure{et} $x_k\handwritten{\unsure{\notin}}X_{f_j}$ \unsure{pour} $k\handwritten{\unsure{\ne}}j$} ; en outre, on sait qu'il y a un entier $q$ tel que chacune des sections $d_i f_j^q$ (où $d_i\handwritten{\otimes}1$ est dans $\underline O_{x_j}$) se prolonge en une section $b_i\handwritten{\in}B$ (I,1,4,th.2). Comme $f_j$ est inversible au point $x_j$ et que les $d_i\handwritten{\otimes}1$ correspondant à $x_j$ sont nulles aux $x_k\handwritten{\ne}x_j$, les éléments $b_i\handwritten{\otimes}1$ forment encore une base de $D\handwritten{\otimes}_C C_y$ sur $C_y$, ce qui établit notre assertion.

Cela étant, \uline{considérons l'homomorphisme} $\handwritten{\Gamma}(Y,\underline O_Y^n)\handwritten{\to}\handwritten{\Gamma}(X,\underline O_X)=B$ \uline{faisant correspondre}
% slanted
\margin{\handwritten{\unsure{ou} \uline{\unsure{bien}} \uline{\unsure{faisant}}}}$b_i$ au i-ème élément de base $(0,..,1,..,\deleted{\unsure{0}}0)$ de $\handwritten{\Gamma}(Y,\underline O_Y^n)$ ; cet homomorphisme définit évidemment un homomorphisme de faisceaux $\boxed{u:\underline A^n\handwritten{\to}\underline B.}$ \deleted{\illegible{1}} Les faisceaux $\underline A^n$ et $\underline B$ sont quasi-cohérents sur $Y$ (II,2,2,prop.5), donc cohérents\inserted{\handwritten{ puisque $Y$ est \unsure{nœthérien|noethérien} et $B$ de type fini par hypothèse}}, et il en est dnnc de même du noyau et du conoyau de $u$. Or, par construction, si \inserted{$s$ est} une section de $\operatorname{Ker}u$ ou de $\operatorname{Coker}u$ au-dessus d'un voisinage de $y$, le germe de $s$ au point $y$ est nul, donc $s$ est nulle dans un voi\deleted{\unsure{s}}sinage $y$ ; on en conclut qu'il existe un voisinage ouvert $V$ de $y$ tel que la restriction $u\handwritten{\vert}V$ soit un \uline{isomorphisme} $\underline A^n\handwritten{\vert}V\handwritten{\to}\underline B\handwritten{\vert}V$.

Cela étant, en vertu du th.2, il suffit de prouver que pour tou\handwritten{t}
% archive sheet number: 3
